% GENERATED FILE. Edit canonical lessons, then run npm run generate:books.
% canonical-source-hash: fnv1a64:a00d78c3ebcf3a95
% canonical-lessons: KA-C231-sahodyogi, KA-C231-sabhe, KA-C231-sandarsana, KA-C231-gumasta, KA-C231-adhikari

\chapter{Colleague, Meeting, Interview, Clerk, Officer}
\label{ch:ka-colleague-meeting-interview-clerk-officer}
\hlchaptermodality{231}

\begin{chapteropening}
\textbf{By the end of this chapter:} \emph{I can say a colleague, a meeting, an interview, a clerk and an officer.}

Complete the last lesson of chapter 231: I can say a colleague, a meeting, an interview, a clerk and an officer.
\end{chapteropening}

\section[\texorpdfstring{\kn{ಸಹೋದ್ಯೋಗಿ} (sahōdyōgi)}{ಸಹೋದ್ಯೋಗಿ (sahōdyōgi)}]{\kn{ಸಹೋದ್ಯೋಗಿ} (sahōdyōgi) — a colleague}

\label{lesson:KA-C231-sahodyogi}



\begin{quote}
Before the new one: say the Kannada for a hoe, then the Kannada for a sickle.
\end{quote}

\subsection*{You'll want to know: \kn{ಸಹೋದ್ಯೋಗಿ}}
\textbf{\kn{ಸಹೋದ್ಯೋಗಿ}} — \emph{sahōdyōgi} — \textquotedblleft{}a colleague\textquotedblright{}.

\begin{cousinweb}
\textbf{\kn{ಸಹೋದ್ಯೋಗಿ}} (\emph{sahōdyōgi}) is Sanskrit \emph{saha}, \textquotedblleft{}with\textquotedblright{}, and \emph{udyogin}, \textquotedblleft{}one who works\textquotedblright{}, from \emph{udyoga}, \textquotedblleft{}exertion, work\textquotedblright{} (CDIAL 2055): \textquotedblleft{}one who works with\textquotedblright{}.
\end{cousinweb}

\begin{grammarlens}[title={What you've built}]
One more word for this chapter.
\end{grammarlens}

\subsection*{Your turn}
\begin{itemize}
\raggedright
  \item \emph{Say it:} \emph{sahōdyōgi}
  \item \emph{Say it:} \emph{sahōdyōgi}, once more
  \item \emph{Say it:} say \emph{sahōdyōgi}
  \item \emph{Recall:} say the Kannada for a hoe, then the Kannada for a sickle, then say \emph{sahōdyōgi} again
\end{itemize}

\begin{tcolorbox}[breakable,colback=teal!4,colframe=teal!35!black,title={Before you move on}]
What does \kn{ಸಹೋದ್ಯೋಗಿ} mean? (\textquotedblleft{}A colleague\textquotedblright{}.) Say it once more.
\end{tcolorbox}
\section[\texorpdfstring{\kn{ಸಭೆ} (sabhe)}{ಸಭೆ (sabhe)}]{\kn{ಸಭೆ} (sabhe) — a meeting}

\label{lesson:KA-C231-sabhe}



\begin{quote}
Before the new one: say the Kannada for a sickle, then the Kannada for a colleague.
\end{quote}

\subsection*{You'll want to know: \kn{ಸಭೆ}}
\textbf{\kn{ಸಭೆ}} — \emph{sabhe} — \textquotedblleft{}a meeting\textquotedblright{}.

\begin{cousinweb}
\textbf{\kn{ಸಭೆ}} (\emph{sabhe}) is Sanskrit \emph{sabhā}, \textquotedblleft{}assembly, hall; council\textquotedblright{}, taken over whole as a learned loan, a tatsama, with Kannada's \emph{-e} for the Sanskrit \emph{-ā} (CDIAL 13167).
\end{cousinweb}

\begin{grammarlens}[title={What you've built}]
One more word for this chapter.
\end{grammarlens}

\subsection*{Your turn}
\begin{itemize}
\raggedright
  \item \emph{Say it:} \emph{sabhe}
  \item \emph{Say it:} \emph{sabhe}, once more
  \item \emph{Say it:} say \emph{sabhe}
  \item \emph{Recall:} say the Kannada for a sickle, then the Kannada for a colleague, then say \emph{sabhe} again
\end{itemize}

\begin{tcolorbox}[breakable,colback=teal!4,colframe=teal!35!black,title={Before you move on}]
What does \kn{ಸಭೆ} mean? (\textquotedblleft{}A meeting\textquotedblright{}.) Say it once more.
\end{tcolorbox}
\section[\texorpdfstring{\kn{ಸಂದರ್ಶನ} (sandarśana)}{ಸಂದರ್ಶನ (sandarśana)}]{\kn{ಸಂದರ್ಶನ} (sandarśana) — an interview}

\label{lesson:KA-C231-sandarsana}



\begin{quote}
Before the new one: say the Kannada for a colleague, then the Kannada for a meeting.
\end{quote}

\subsection*{You'll want to know: \kn{ಸಂದರ್ಶನ}}
\textbf{\kn{ಸಂದರ್ಶನ}} — \emph{sandarśana} — \textquotedblleft{}an interview\textquotedblright{}.

\begin{cousinweb}
\textbf{\kn{ಸಂದರ್ಶನ}} (\emph{sandarśana}) is Sanskrit \emph{sandarśana}, \textquotedblleft{}view; seeing, showing\textquotedblright{}, taken over whole as a learned loan, a tatsama (CDIAL 12899). A seeing face to face became an interview.
\end{cousinweb}

\begin{grammarlens}[title={What you've built}]
One more word for this chapter.
\end{grammarlens}

\subsection*{Your turn}
\begin{itemize}
\raggedright
  \item \emph{Say it:} \emph{sandarśana}
  \item \emph{Say it:} \emph{sandarśana}, once more
  \item \emph{Say it:} say \emph{sandarśana}
  \item \emph{Recall:} say the Kannada for a colleague, then the Kannada for a meeting, then say \emph{sandarśana} again
\end{itemize}

\begin{tcolorbox}[breakable,colback=teal!4,colframe=teal!35!black,title={Before you move on}]
What does \kn{ಸಂದರ್ಶನ} mean? (\textquotedblleft{}An interview\textquotedblright{}.) Say it once more.
\end{tcolorbox}
\section[\texorpdfstring{\kn{ಗುಮಾಸ್ತ} (gumāsta)}{ಗುಮಾಸ್ತ (gumāsta)}]{\kn{ಗುಮಾಸ್ತ} (gumāsta) — a clerk}

\label{lesson:KA-C231-gumasta}



\begin{quote}
Before the new one: say the Kannada for a meeting, then the Kannada for an interview.
\end{quote}

\subsection*{You'll want to know: \kn{ಗುಮಾಸ್ತ}}
\textbf{\kn{ಗುಮಾಸ್ತ}} — \emph{gumāsta} — \textquotedblleft{}a clerk\textquotedblright{}.

\begin{cousinweb}
\textbf{\kn{ಗುಮಾಸ್ತ}} (\emph{gumāsta}) is a loan from Persian \emph{gumāšta}, \textquotedblleft{}an appointed agent\textquotedblright{}, which came through Hindustani.
\end{cousinweb}

\begin{grammarlens}[title={What you've built}]
One more word for this chapter.
\end{grammarlens}

\subsection*{Your turn}
\begin{itemize}
\raggedright
  \item \emph{Say it:} \emph{gumāsta}
  \item \emph{Say it:} \emph{gumāsta}, once more
  \item \emph{Say it:} say \emph{gumāsta}
  \item \emph{Recall:} say the Kannada for a meeting, then the Kannada for an interview, then say \emph{gumāsta} again
\end{itemize}

\begin{tcolorbox}[breakable,colback=teal!4,colframe=teal!35!black,title={Before you move on}]
What does \kn{ಗುಮಾಸ್ತ} mean? (\textquotedblleft{}A clerk\textquotedblright{}.) Say it once more.
\end{tcolorbox}
\section[\texorpdfstring{\kn{ಅಧಿಕಾರಿ} (adhikāri)}{ಅಧಿಕಾರಿ (adhikāri)}]{\kn{ಅಧಿಕಾರಿ} (adhikāri) — an officer, an official}

\label{lesson:KA-C231-adhikari}



\begin{quote}
Before the new one: say the Kannada for an interview, then the Kannada for a clerk.
\end{quote}

\subsection*{You'll want to know: \kn{ಅಧಿಕಾರಿ}}
\textbf{\kn{ಅಧಿಕಾರಿ}} — \emph{adhikāri} — \textquotedblleft{}an officer, an official\textquotedblright{}.

\begin{cousinweb}
\textbf{\kn{ಅಧಿಕಾರಿ}} (\emph{adhikāri}) is Sanskrit \emph{adhikārin}, \textquotedblleft{}one with authority, one in charge\textquotedblright{}, from \emph{adhikāra}, \textquotedblleft{}authority\textquotedblright{}. Kannada drops the final \emph{-n}.
\end{cousinweb}

\begin{grammarlens}[title={What you've built}]
One more word for this chapter.
\end{grammarlens}

\subsection*{Your turn}
\begin{itemize}
\raggedright
  \item \emph{Say it:} \emph{adhikāri}
  \item \emph{Say it:} \emph{adhikāri}, once more
  \item \emph{Say it:} say \emph{adhikāri}
  \item \emph{Recall:} say the Kannada for an interview, then the Kannada for a clerk, then say \emph{adhikāri} again
\end{itemize}

\begin{tcolorbox}[breakable,colback=teal!4,colframe=teal!35!black,title={Before you move on}]
What does \kn{ಅಧಿಕಾರಿ} mean? (\textquotedblleft{}An officer, an official\textquotedblright{}.) Say it once more.
\end{tcolorbox}
